%% --------------------------------------------------------
\section{Раманівський спектр}
%% --------------------------------------------------------

\textbf{Раман-спектроскопія} досліджує не поглинання, а \emph{розсіювання} світла на молекулах.
Коли фотон із частотою $\nu_0$ падає на молекулу, більшість випромінюється назад із тією ж частотою — це \emph{Релеївське розсіювання}.
Однак невелика частина розсіюється з іншою частотою, зсуненою на коливальну величину $\nu_k$:
\[
\nu_{\text{розс}} = \nu_0 \pm \nu_k,
\]
де «мінус» відповідає \emph{лінії Стокса}, а «плюс» — \emph{анти-Стоксу}.

\textbf{Фізичний зміст.}
Світлове поле $\mathbf{E}$ індукує у молекулі миттєвий дипольний момент:
\[
\boldsymbol{\mu}_{\text{інд}} = \boldsymbol{\alpha}\,\mathbf{E},
\]
де $\boldsymbol{\alpha}$ — тензор поляризовності, що описує, як електронна хмара реагує на поле.
Під час коливання молекули параметр $\boldsymbol{\alpha}$ змінюється.
Саме ця зміна створює компоненту поля на частоті $\nu_0 \pm \nu_k$ — тобто Раманівське розсіювання.

\textbf{Умова активності.}
Коливальна мода $Q_k$ є Раман-активною, якщо при русі ядер поляризовність змінюється:
\[
\frac{\partial \boldsymbol{\alpha}}{\partial Q_k} \neq 0.
\]
Інтенсивність розсіювання для моди $k$ пропорційна квадрату цієї похідної:
\[
I_k^{\text{Раман}} \propto
\left|
\frac{\partial \alpha_{ij}}{\partial Q_k}
\right|^2,
\]
де індекси $i,j=x,y,z$ позначають компоненти тензора поляризовності.

\vspace{1em}
\textbf{Як моделюють Раман-спектр на практиці:}
\begin{enumerate}
    \item \textbf{Коливальні моди.}
    З оптимізованої геометрії обчислюють гесіан $\mathbf{H}$ та отримують нормальні частоти $\omega_k$ і вектори мод $\mathbf{L}_k$.

    \item \textbf{Симетрія.}
    За допомогою теорії груп визначають, які моди Раман-активні (для цього аналізують перетворення тензора $\boldsymbol{\alpha}$ у даній точковій групі).

    Для молекул із центром інверсії (наприклад, \ce{CO2}) діє \emph{правило взаємовиключності}:
    \begin{itemize}
        \item парні моди ($g$) — \textbf{Раман-активні}, але ІЧ-неактивні;
        \item непарні моди ($u$) — \textbf{ІЧ-активні}, але Раман-неактивні.
    \end{itemize}

    \item \textbf{Похідна поляризовності.}
    Для кожної моди обчислюють зміну тензора $\boldsymbol{\alpha}$ при малих зміщеннях $\pm\Delta Q_k$ вздовж нормального вектора:
    \[
    \left. \frac{\partial \alpha_{ij}}{\partial Q_k} \right|_0
    \approx
    \frac{\alpha_{ij}(+\Delta Q_k) - \alpha_{ij}(-\Delta Q_k)}{2\Delta Q_k}.
    \]

    \item \textbf{Інтенсивність лінії.}
    Вводять два інваріанти тензора похідних:
    \[
    \bar{\alpha}_k = \tfrac{1}{3} \mathrm{tr}\!\left(\frac{\partial \boldsymbol{\alpha}}{\partial Q_k}\right),
    \qquad
    \gamma_k^2 = \tfrac{1}{2}\!\!\sum_{i<j}\!\!\left[\!
    \left(\frac{\partial \alpha_{ii}}{\partial Q_k} - \frac{\partial \alpha_{jj}}{\partial Q_k}\right)^{\!2}
    + 6\!\left(\frac{\partial \alpha_{ij}}{\partial Q_k}\right)^{\!2}\right]\!.
    \]
    Тоді відносна Раман-активність:
    \[
    I_k = 45\,\bar{\alpha}_k^2 + 7\,\gamma_k^2.
    \]

    \item \textbf{Побудова спектра.}
    Кожній моді з частотою $\nu_k$ приписують інтенсивність $I_k$ та накладають профіль лінії (зазвичай Лоренців, ширина $\Gamma \approx 5~\text{см}^{-1}$):
    \[
    I(\nu) = \sum_k I_k \, g(\nu - \nu_k).
    \]
\end{enumerate}

\textbf{Резюме.}
Якщо при коливанні молекули змінюється її поляризовність — мода видно в Раман-спектрі.
Якщо ж змінюється дипольний момент — мода проявляється в ІЧ-спектрі.
Тому ІЧ та Раман часто дають взаємодоповнюючу інформацію про коливальні стани молекули.

%% --------------------------------------------------------
\subsubsection{Раманівського спектр молекули \ce{H2O}}
%% --------------------------------------------------------


В програмі \inputcode{h2o_raman_activity.py} наведено розрахунок раманівського спектру молекули \ce{H2O}.

\textbf{Типові активності для \ce{H2O}:}
\begin{center}
    \begin{tabular}{lcc}
        \toprule
        Мода    & $\nu$ (см$^{-1}$) & Раман-активність (Å$^4$/amu) \\
        \midrule
        $\nu_1$ & 3657              & 1.8                          \\
        $\nu_2$ & 1595              & 0.4                          \\
        $\nu_3$ & 3756              & 3.2                          \\
        \bottomrule
    \end{tabular}
\end{center}


\textbf{Принцип взаємного виключення:}

Для молекул з центром інверсії (наприклад, \ce{CO2}):
\begin{itemize}
    \item Парні моди ($g$): Раман активні, ІЧ неактивні
    \item Непарні моди ($u$): ІЧ активні, Раман неактивні
\end{itemize}

%\subsubsection{Раман спектр бензену}
%
%\inputcode{benzene_raman_spectrum.py}
%
%\textbf{Вибрані моди бензену \ce{C6H6} (D$_{6h}$):}
%
%\begin{center}
%    \begin{tblr}{
%        colspec = {X[c] X[c] X[r] X[c] X[c]},
%        row{1} = {font=\bfseries},
%        hline{1,2,Z} = {solid},
%            }
%        Мода       & Симетрія & Частота     & ІЧ        & Раман            \\
%                   &          & (см$^{-1}$) &                              \\
%        $\nu_1$    & A$_{1g}$ & 993         & Неактивна & Активна          \\
%        $\nu_2$    & A$_{1g}$ & 3062        & Неактивна & Активна (сильна) \\
%        $\nu_6$    & E$_{2g}$ & 608         & Неактивна & Активна          \\
%        $\nu_{18}$ & E$_{1u}$ & 1038        & Активна   & Неактивна        \\
%        $\nu_{19}$ & E$_{1u}$ & 3080        & Активна   & Неактивна        \\
%    \end{tblr}
%\end{center}
%
%\textit{Розрахунок: B3LYP/6-311+G(2d,p)}
%
%%---------------------------------------------------------
%\begin{figure}[h!]\centering
%%    \includegraphics[width=0.85\linewidth]{\currfiledir/benzene_raman_spectrum.pdf}
%    \caption{Раман спектр бензену демонструє принцип взаємного виключення.
%        Симетричні моди ($g$) активні в Раман, але неактивні в ІЧ.}
%    \label{pic:benzene_raman}
%\end{figure}
%%---------------------------------------------------------
%
%\textbf{Характерні Раман смуги:}
%\begin{itemize}
%    \item 993 см$^{-1}$: дихальна мода кільця (ring breathing)
%    \item 3062 см$^{-1}$: симетричне валентне \ce{C-H}
%    \item 608 см$^{-1}$: деформація кільця
%\end{itemize}
%
%\subsubsection{Порівняння ІЧ та Раман}
%
%\begin{center}
%    \begin{tblr}{
%        colspec = {X[l,2] X[l,3] X[l,3]},
%        row{1} = {font=\bfseries},
%        hline{1,2,Z} = {solid},
%            }
%        Властивість               & ІЧ спектроскопія         & Раман спектроскопія         \\
%        Правило відбору           & $\frac{d\mu}{dQ} \neq 0$ & $\frac{d\alpha}{dQ} \neq 0$ \\
%        Джерело                   & ІЧ лампа                 & Лазер (видиме/УФ)           \\
%        Зразок                    & Розчин, плівка, KBr      & Розчин, кристал             \\
%        Вода                      & Сильно поглинає          & Слабкий сигнал              \\
%        Скло                      & Непрозоре                & Прозоре                     \\
%        Полярні групи             & Сильний сигнал           & Слабкий                     \\
%        Неполярні                 & Слабкий                  & Сильний сигнал              \\
%        \ce{C=C}, \ce{C\bond{3}C} & Слабкі/неактивні         & Сильні                      \\
%        \ce{C=O}, \ce{O-H}        & Дуже сильні              & Слабкі                      \\
%    \end{tblr}
%\end{center}